\section{Methodology and Adaptation}
\label{sec:method}

\TODO{Give enough detail that a reader could reimplement the model without the code.}

\begin{figure}[t]
  \centering
  \placeholderfig{2.5in}
  \caption{Architecture overview. \TODO{Mark which components come from the source paper and which are new.}}
  \label{fig:arch}
\end{figure}

\subsection{Baseline Architecture}
\label{sec:method-baseline}
\TODO{The unmodified source architecture, with notation and equations.}

\subsection{Structural Modification}
\label{sec:method-mod}
\TODO{The new module, loss term, or embedding, with equations (\cref{fig:arch}).}

\subsection{Adaptation to the New Dataset}
\label{sec:method-adapt}
\TODO{Every change needed to run the original code on the new data: input layer changes, loss function adjustments, data loader rewrites.}

\subsection{Theoretical Rationale}
\label{sec:method-rationale}
\TODO{Why each variation or substitution should help on this data.}
